\section{Problem formulation}\label{sec:problem}
A central planner is called whenever a replanning trigger fires. Each UAV keeps the task that it is currently flying to
or serving, which determines when and where it becomes available again and how much energy it has left. Each open task
$j$ has a location, a release time $r_j$, a deadline $d_j$ by which its service must be completed, a service duration
and a reward $w_j$. A plan assigns an ordered sequence of open tasks to each UAV, and tasks may remain unassigned, so
the planning problem is a prize-collecting routing problem with time
windows~\cite{solomon1987algorithms,vansteenwegen2011orienteering}. Arrival, start and finish times along a route are
computed as usual for time windows~\cite{savelsbergh1992vrptw}. A route is feasible if every service ends before its
deadline and the UAV returns to the depot with an energy reserve of 10\% of its capacity. A plan $\pi$ is evaluated by
\begin{equation}\label{eq:obj}
J(\pi)=\alpha R(\pi)-\gamma E(\pi)-\delta C(\pi),
\end{equation}
where $R$ is the reward of the tasks served on time, $E$ the flown distance in kilometres and $C$ the number of tasks
that move to another UAV than in the previous plan, with $\alpha=1$, $\gamma=0.05$ per kilometre and $\delta=0.3$. $W$
denotes the total reward of the open tasks.

Computation is not free. A plan computed for $\tau$ milliseconds is released $\kappa\tau$ milliseconds later in mission
time, where the latency factor $\kappa$ represents planning computers slower than the evaluation computer. An idle UAV
waits for the new plan, so its route starts later, and tasks whose deadline can no longer be met are lost. The plan
then reaches the vehicles after a dissemination latency $L$ that the planner does not control.

The planner uses a nominal model with straight legs at the cruise speed of each UAV and nominal service durations. The
executed finish time of task $j$ is $F_j=f_j+e_j$, where $f_j$ is the nominal finish time at publication and $e_j$ an
execution error caused by acceleration, wind, speed tracking and service-time variability. A task yields its reward only
if $F_j\le d_j$. As in dispatching practice, a UAV skips a task at departure if the nominal estimate from its actual
state already misses the deadline, and the task returns to the pool. The utility of a mission is the reward of all tasks
served on time divided by the total reward of the mission. Three replanning triggers are considered. The event trigger
replans on task arrivals, UAV failures and deadline changes, the periodic trigger in addition every $T$ seconds, and
the deviation trigger in addition when a finish time deviates from the plan by more than $\theta$ seconds. During
planning, every UAV finishes its committed task and then waits, and events are queued.
